成熟组件进入购买系统以后，组合方式取决于业务事实的归属。前五章分别展开复制组、协调服务、Redis、分布式存储、事务和执行运行时，本节把它们放回两套完整部署中，观察一份购买怎样经过这些组件。两套部署选择不同的业务权威实现，外部资金、消息、消费和履约继续沿持久身份连接。

\subsection{分区状态机承载购买权威}

第一套部署采用本书的完整分区 Raft 状态机。入口按账户找到 b1，G1 在一条受理命令中核验会话、仲裁完整 $S$、固定输入并创建 P/C/T。需要组合库存时，协调责任与各库存组通过原 TX 形成决定。时间、随机数和平台 HTTP 在状态机外取得，状态机应用使用已经确定的命令输入。工作者领取责任以后调用外部平台，确定结果通过新的权威命令回填。

\CTwentyOneFigure{ch21-12}{工业组件围绕唯一业务权威组合，缓存、消息、执行与归档分别承担自己的职责}{fig:ch21-components}

控制面可采用第十七章介绍的配置服务。其归属和迁移记录决定 b1 当前承载在哪个组，工作代次在相应资源端核验。Redis 保存查询投影、路由缓存或可重建索引，命中可以降低观察延迟，miss 则读取原权威。消息域通过适配器一次选择 Kafka 或 RabbitMQ，发布器从权威 Outbox 读取完整 EV，消费者在自己的事务中保存 I/B/F 与 dispatch 意图。

取 Alice 的 3000 分购买为例：入口从 G1 得到 \code{P:b1:1/C:b1:1/T:b1:1}。以下把对应的 \code{O:b1:1} 和 \code{EV:b1:1} 局部简写为 O1、EV1。平台通过原范围和键保存 X1。G1 接受 X1 后，把原 provider result、payment result、O1 和 EV1 一同固定。发布器先把 EV1 写入 Kafka 分区 offset 80，随后确认响应丢失，恢复后又写到 offset 81。消费者在 offset 80 对 EV1 保存 I 与 F501，offset 81 读取到原 I 和业务守卫，沿原 F501 继续。dispatch 重启按原身份查询发运对象，将查询所得结果回填到 F501。

改选 RabbitMQ 时，EV1 经 exchange 路由到 quorum queue，发布确认和手工 ACK 分别连接发布器与消费者。消费事务提交后 ACK 丢失会产生重投，消费者读取原 I/B/F 继续处理。适配器负责消息交付和进度管理，G1 保存购买结果，消费域保存履约事实。部署把 broker 副本、消息磁盘、保留策略、消费并发与恢复速率纳入前两节的同一预算。

归档层用 HDFS 或对象存储保存不可变图像、账本导出与固定批量输入。MapReduce、Spark 读取固定来源，生成候选输出并通过提交协议发布。Ray 负责动态任务、对象和 Actor，函数执行器负责短生命周期调用。它们执行的修复动作从原 C/T 或 TX 取得身份与资格。新运行时增加自己的执行历史、调度与恢复容量，购买责任继续从权威枚举。

这套部署适合业务状态机边界清楚、按分区扩展且愿意维护命令和快照协议的系统。所有业务相关可变状态共同进入复制与迁移，快照模式升级、状态机确定性、日志增长和复制运维成为长期成本。例如新增退款时，命令要核验原收款与累计占额，图像要保存 RF 及平台证据，恢复器要枚举未知退款，独立检查器则核对成功退款与未知占额之和。新增业务状态由此进入同一份复制和恢复历史。

\subsection{分布式 SQL 承载另一种权威实现}

第二套部署以分布式 SQL 数据库承载购买权威。入口、恢复服务和消费者仍是独立进程，S 绑定、原输入、C/T、结果、订单和 Outbox 映射为表与约束。数据库提供行、索引和跨分片事务，应用用事务保持这些关系。TiDB/TiKV 的 Region 与 Peer 结构见 17.5 节，其跨 Region 事务以及 CockroachDB 的提交与恢复路径见第十九章。这些数据库各自的复制、时间与提交机制在数据库内部工作。

\CTwentyOneFigure{ch21-21}{分布式 SQL 方案把业务原子写入交给数据库事务，外部效果沿原责任和独立接口恢复}{fig:ch21-sql-composition}

以外部付款购买为例，第一段事务插入完整 S 唯一绑定及原输入，生成实际 P/C/T，并保存付款状态。遇到唯一冲突时读取已存在的原绑定，比较完整语义输入。事务提交返回的实际编号继续关联后台任务。网络响应未知时，入口按原 S 查询，读取到的绑定决定究竟恢复哪份购买。

付款工作者用短事务领取责任、取得工作代次并读取保存的原键，事务结束以后发送平台 HTTP。X1 返回以后，再用事务核验代次、原输入与当前结果，固定原结果、O 和 EV。若发送期间一直持有数据库事务，等待平台的时间会占用连接与锁；平台提交 X1 后，即使本地事务回滚，收款也已经发生。因此领取与回填各用一段短事务，持久 C/T 连接两段事务之间的平台调用与查询。

购买表和库存表位于同一分布式数据库、由同一事务服务管理时，应用可以用一份跨分片事务更新它们，数据库内部协议负责准备、提交和恢复。若库存属于独立服务或另一数据库，购买协调者创建 TX，各库存服务按 TX 保存准备占用，协调者再持久保存提交或中止决定。恢复者读取这份决定，逐一通知参与者完成占用的兑现或释放。部署中的事务服务与数据归属，决定应用采用哪一条提交路径。

数据库隔离级别还影响应用中的读取和重试。TiDB/TiKV 7.5.8 与 CockroachDB 25.4.0 分别采用第十九章给定的配置、事务模式与错误处理路径。序列化重试重新执行事务体时，原 S 和输入保持，外部 HTTP 在事务结束以后发送。唯一冲突触发原绑定读取，实际提交的编号连接后台责任；领取或回填影响零行时，工作者读取当前代次与已有结果，沿原责任决定下一步。

这套部署减少自定义复制、SQL 查询与在线数据维护的负担，代价转向索引写放大、热点事务、连接池、数据库版本与模式升级。把全部记录放在一张表上，可能造成单索引热点；把相关记录分散太远，又会增加事务参与分片。数据布局依据主要约束和查询路径选择，配套容量压测同时覆盖新受理、同身份恢复与故障后回填。

\subsection{组件角色、版本和授权共同进入部署清单}

工业部署中，同一种组件会进入多条业务链。Kafka 官方使用页面介绍，LinkedIn 用它承载活动流和运行指标，PayPal 的用途包括应用指标、数据库同步、日志聚合与风险分析，Uber 将其用于在线及近实时场景，Wikimedia 则用它连接生产事件与分析，包括 Wikipedia 缓存失效和向 Hadoop 接入数据流。PingCAP 的客户材料介绍，Square 采用 TiDB 扩展增长中的 MySQL 工作负载，Bolt 则围绕 AWS 上大量微服务的数据需求推进 MySQL 架构的扩展~\cite{ch21-kafka-adoption,ch21-tidb-adoption}。

这些用途对应着不同的组件角色。活动流、指标和风险分析需要把同一批事件交给多个处理者，Kafka 的分区日志让各消费组独立读取历史。购买系统可据此把 EV1 分别交给履约、风控与分析，业务消费者各保存自己的 I/B。Redis 把可重建观察移到低延迟路径，权威承担首次受理和原结果。分布式 SQL 则承载持续增长的业务表与多表约束，应用把外部付款接到持久 C/T。日志保留、缓存重建和数据库恢复分别占用磁盘、网络与值守能力，组件的部署预算沿这些实际职责分配。

本书组合使用 Redis 8.10.2、TiDB/TiKV 7.5.8、CockroachDB 25.4.0、Hadoop 3.4.2、Spark 4.1.3、Ray 2.59.0，以及第十六章的消息组件版本。版本选择连接维护周期、数据格式和兼容要求。升级后的代码按原模式读取已有 S 和原结果，同 S 异参返回冲突，未决 TX 与旧工作流从保存位置继续。流量逐步转入新实例，接管者按这些记录恢复尚未完成的责任。

版本固定以后，部署清单还要写明组件是供内部业务使用、分发修改版，还是向第三方提供托管功能。软件许可分别约束这些使用方式。Hadoop、Spark、Ray、Kafka 等相关组件采用 Apache 2.0，它允许按条款使用、修改和分发，提供相应专利授权，分发时保留许可、相关通知并标注修改。团队自行维护修改版时，把升级、支持和故障处置的人力一同列入运行成本。

Redis 8 的三种可选许可把上述使用方式连到不同条件。源码可用许可 RSALv2 限制把软件功能提供给第三方作为服务，具体范围由其服务和产品定义判断。服务端公共许可 SSPLv1 对这类服务要求开放服务源码，范围还包含用于提供服务的管理、监控、备份和托管等程序。GNU Affero 通用公共许可 AGPLv3 则要求修改版向远程交互用户提供对应源码，并另有分发条件。内部购买系统的缓存、向客户开放 Redis 功能的托管平台，以及分发修改版，因而分别记录所选许可、部署方式和交付材料。

CockroachDB 25.4.0 的 CockroachDB Software License 是厂商定义的软件许可，除条款界定的有限非生产用途外，运行需要有效许可密钥，并按密钥与许可要求处理使用限制和向厂商报告使用数据的遥测。生产部署据此安排许可密钥、遥测出口和商业支持，分发修改版则核对相应分发条件。部署清单保存实际版本、选用的许可及对应文本，扩容、升级或改变服务方式时，据此调整授权和预算~\cite{ch21-licenses}。

\subsection{上线、升级和恢复是一条持续链路}

两套部署都先给服务和保存域固定版本，再建立持久卷、备份和权限。业务入口、管理入口、平台出口与消息路径分别验证，先只读观察再启用外部发送。配置发布记录谁采用了哪个版本，实例替换时从持久责任恢复，系统可以回答任何一份 C 目前由谁负责。

模式升级先保持旧身份可解释。\code{purchase.create.v1} 使用本地扣款，恢复时读取原输入、本地扣款事实和原结果。\code{purchase.create.external.v2} 引入外部付款，已经受理的身份从记录读取原 provider、scope、平台键、金额和响应语义。报价更新只参与新身份的首次业务判断。新增字段采用兼容的读取和写入顺序，运行中的恢复器按动作与模式版本处理历史。

滚动升级期间，新旧应用实例可能同时存在。外部付款沿原 C 和原键，权威代次拒绝过时回填，工作流部署则为未完成执行保留匹配代码。例如旧实例已受理 \code{external.v2} 付款，平台提交 X1 后进程退出，新实例接管时先取得原 C，读取保存的平台范围和键并查询 X1，再按该记录的结果模式回填。备份定期在独立环境加载，恢复器从图像重建同一关联，继续处理尚未完成的 TX、付款和履约责任。

新增组件时，一份购买会多出相应的保存位置和恢复步骤。例如工作流引擎接替扫描器的控制流程以后，部署增加历史数据库，升级保留能解释未完成历史的代码，值守者从工作流定位原 C，再从业务权威读取付款结果。对象存储承载冷状态与图像时，上传者先保存候选对象，随后发布带来源位置和摘要的有效图像，恢复者固定该图像再回放后续日志。历史保留、兼容代码、图像发布与日志裁剪由此成为具体的维护工作。
